\chapter{Radiative Forcing and Climate Response}\label{chap:ch14_forcing_response}

% --- local symbols (minimal, chapter-scoped) ---------------------------
\providecommand{\Forc}{F}                 % radiative forcing [W/m^2]
\providecommand{\Qflux}{\overline{Q}}     % frozen q-flux (time-mean restoring flux)
\providecommand{\Hrest}{\lambda}          % Haney restoring coefficient [W/m^2/K]
\providecommand{\SST}{T_{\mathrm{s}}}     % sea-surface temperature
\providecommand{\SSTt}{T^{\star}}         % SST restoring target (WOA climatology)

A climate model earns its name only if it \emph{responds}: if you push it with an
external forcing, the simulated climate must move, and move by an amount you can
defend. This chapter is about giving \model{} that capability for the single most
important anthropogenic lever, atmospheric carbon dioxide. We proceed in three
steps that mirror the code. First, we convert a CO\textsubscript{2} concentration
into a radiative forcing through the Myhre logarithmic law
(\S\ref{sec:f14_myhre}). Second, we confront a structural obstacle: the surface
flux correction that keeps the control climate from drifting also \emph{absorbs}
any forcing we impose, so we must free the ocean surface with the q-flux method
(\S\ref{sec:f14_qflux}). Third, we design the abrupt-$2\times$CO\textsubscript{2}
experiment as a difference of two branches (\S\ref{sec:f14_experiment}) and read
off the response (\S\ref{sec:f14_result}). Throughout, the central honesty of the
chapter is that what \model{} produces is a \emph{transient, surface-forcing
proxy} for the transient climate response (TCR), not the equilibrium climate
sensitivity (ECS); \S\ref{sec:f14_honesty} makes that distinction quantitative.

This chapter builds directly on the coupler and surface-flux budget of
Chapter~\ref{chap:ch13_coupler_fluxes}; the ocean that does the warming is the
prognostic core of Chapter~\ref{chap:ch07_ocean_prognostic_core}, and the
overturning circulation that the same forcing can perturb is the subject of
Chapter~\ref{chap:ch17_amoc_tipping}.

% ======================================================================
\section{The Myhre CO\texorpdfstring{$_2$}{2} radiative forcing}\label{sec:f14_myhre}
% ======================================================================

The radiative effect of a well-mixed greenhouse gas is, to leading order, not
linear in its concentration but logarithmic: each \emph{doubling} of
CO\textsubscript{2} adds roughly the same forcing because the strong absorption
lines saturate and only the band wings contribute further. \citet{myhre1998}
fitted line-by-line radiative-transfer calculations to the compact expression

\begin{equation}
  \Forc(C) \;=\; 5.35 \,\ln\!\left(\frac{C}{C_0}\right)
  \qquad [\mathrm{W\,m^{-2}}],
  \label{eq:f14_myhre}
\end{equation}

where $C$ is the CO\textsubscript{2} mixing ratio in parts per million by volume
(ppm), $C_0=280\,\mathrm{ppm}$ is the pre-industrial reference, and the
coefficient $5.35\,\mathrm{W\,m^{-2}}$ carries the radiative physics. By
construction $\Forc(C_0)=0$, and a doubling $C=2C_0$ gives the canonical
\begin{equation}
  \Forc(2C_0) \;=\; 5.35\,\ln 2 \;\approx\; 3.71\,\mathrm{W\,m^{-2}}.
  \label{eq:f14_2x}
\end{equation}

\model{} implements Equation~\eqref{eq:f14_myhre} verbatim as a pure-\jax{}
function, so it is differentiable in $C$ --- a property we will exploit when we
ask for $\mathrm{d}(\text{climate})/\mathrm{d}(\text{CO}_2)$ in
Chapter~\ref{chap:ch15_differentiable_applications}.

\begin{lstlisting}[caption={Myhre forcing, the single radiative channel.}]
def co2_forcing_wm2(co2_ppm, co2_ref=280.0):
    return 5.35 * jnp.log(jnp.asarray(co2_ppm) / co2_ref)
\end{lstlisting}
\codref{chronos\_esm/coupler/dino\_coupling.py}{\code{co2\_forcing\_wm2} ---
Equation~\eqref{eq:f14_myhre}}

\begin{remark}[A single forcing channel, and where it enters]
\model{} injects $\Forc$ at exactly one point: it is \emph{added to the ocean
surface heat budget} as an extra downwelling term, alongside the shortwave,
longwave, sensible and latent fluxes assembled in
Chapter~\ref{chap:ch13_coupler_fluxes}. Writing the net downward surface heat
flux of the control as $Q_{\mathrm{net}}$, the forced budget is simply
\begin{equation}
  Q_{\mathrm{net}}^{\,\mathrm{forced}} \;=\; Q_{\mathrm{net}} \;+\; \Forc(C).
  \label{eq:f14_inject}
\end{equation}
A tempting alternative --- shifting the SST-slaved atmospheric radiative
equilibrium temperature $T_{\mathrm{eq}}$ --- would be nearly inert (the
single-/multi-level atmosphere relaxes back) and would risk double-counting, so
it is deliberately \emph{not} used. The code comment in
\code{co2\_forcing\_wm2} states this design choice explicitly.
\end{remark}

The forcing is added in \code{ocean\_fluxes\_jax} \emph{before} the
flux-correction term, because --- as we now show --- that correction is precisely
the artifact that would otherwise swallow it.

\begin{lstlisting}[caption={Forcing enters the surface budget before flux correction.}]
net_heat = sw_net + lw_down - lw_up - sens - lat
if co2_ppm is not None:
    net_heat = net_heat + co2_forcing_wm2(co2_ppm)   # single channel
if sst_target is not None:
    net_heat = net_heat + flux_correction.heat_correction(
        sst_K, sst_target, restore_tau_days, q_flux=q_flux)
\end{lstlisting}
\codref{chronos\_esm/coupler/dino\_coupling.py}{\code{ocean\_fluxes\_jax} ---
forcing injected before the SST flux correction}

% ======================================================================
\section{The flux-correction obstacle and the q-flux cure}\label{sec:f14_qflux}
% ======================================================================

\subsection{Why strong restoring eats the forcing}

The control climate of \model{} is held against drift by a \emph{Haney
restoring} of SST toward the WOA18 climatology $\SSTt$. The correction is a
linear damping of the SST error,
\begin{equation}
  Q_{\mathrm{corr}}(\SST) \;=\; \Hrest\,\bigl(\SSTt - \SST\bigr),
  \qquad
  \Hrest \;=\; \frac{\densz\, c_p\, \Delta z_1}{\tau},
  \label{eq:f14_haney}
\end{equation}
where $\Hrest$ is the restoring coefficient $[\mathrm{W\,m^{-2}\,K^{-1}}]$,
$\densz=1025\,\mathrm{kg\,m^{-3}}$ and $c_p=3985\,\mathrm{J\,kg^{-1}K^{-1}}$ are
seawater density and heat capacity, $\Delta z_1\approx 50\,\mathrm{m}$ is the
thickness of the top ocean layer, and $\tau$ is the restoring timescale. The
control uses $\tau=30\,\mathrm{d}$, which gives a \emph{strong}
\begin{equation}
  \Hrest \;=\; \frac{1025\times 3985\times 50}{30\times 86400}
  \;\approx\; 79\,\mathrm{W\,m^{-2}\,K^{-1}}.
  \label{eq:f14_lambda30}
\end{equation}

\codref{chronos\_esm/ocean/flux\_correction.py}{\code{restoring\_lambda},
\code{restoring\_flux} --- Equations~\eqref{eq:f14_haney}--\eqref{eq:f14_lambda30}}

Now superimpose a forcing $\Forc$ on this surface budget and ask for the
\emph{steady} SST anomaly $\Delta\SST$ it can sustain. The surface heat balance
of the perturbed top layer requires the extra input $\Forc$ to be removed by the
restoring it provokes,
\begin{equation}
  \Forc \;-\; \Hrest\,\Delta\SST \;=\; 0
  \quad\Longrightarrow\quad
  \Delta\SST \;=\; \frac{\Forc}{\Hrest}.
  \label{eq:f14_balance}
\end{equation}
With $\Forc=3.71\,\mathrm{W\,m^{-2}}$ and $\Hrest\approx
79\,\mathrm{W\,m^{-2}\,K^{-1}}$, Equation~\eqref{eq:f14_balance} permits only
$\Delta\SST\approx 0.047\,\mathrm{K}$ --- a forty-seventh of a degree. The strong
restoring is, in effect, an infinite heat reservoir clamped to the present
climate: it silently supplies whatever flux is needed to hold $\SST\to\SSTt$ and
therefore \emph{absorbs the forcing entirely}. This is the central problem of
forced experiments with a flux-corrected model.

\subsection{The q-flux methodology}

The cure is the classical mixed-layer \emph{q-flux} (flux-anomaly correction).
The key realization is that the time-mean of the restoring flux over a converged
control \emph{is itself the flux correction the model needs} to hold its mean
state:
\begin{equation}
  \Qflux(\lat,\lon) \;\equiv\;
  \Bigl\langle\, \Hrest\,\bigl(\SSTt - \SST\bigr) \,\Bigr\rangle_t,
  \label{eq:f14_qbar}
\end{equation}
the angle brackets denoting a time average over the equilibrated control. We
\emph{freeze} this field. In the forced (``free'') run we then prescribe the
fixed $\Qflux$ and replace the strong restoring with only a \emph{weak,
long-timescale anomaly} restoring:
\begin{equation}
  Q_{\mathrm{corr}}^{\,\mathrm{free}}(\SST)
  \;=\; \underbrace{\Qflux}_{\text{frozen, holds mean}}
  \;+\; \underbrace{\Hrest_{\mathrm{weak}}\,\bigl(\SSTt - \SST\bigr)}_{\text{bounds drift}},
  \qquad
  \Hrest_{\mathrm{weak}} = \frac{\densz c_p \Delta z_1}{\tau_{\mathrm{weak}}}.
  \label{eq:f14_free}
\end{equation}
With $\tau_{\mathrm{weak}}=3650\,\mathrm{d}$ (ten years) the weak coefficient
collapses to
\begin{equation}
  \Hrest_{\mathrm{weak}}
  \;=\; \frac{1025\times 3985\times 50}{3650\times 86400}
  \;\approx\; 0.65\,\mathrm{W\,m^{-2}\,K^{-1}},
  \label{eq:f14_lambdaweak}
\end{equation}
a factor of $\sim\!120$ weaker than the control restoring. The division of labour
is exact: the frozen $\Qflux$ carries the mean state (so the cold-tongue drift
documented in Chapter~\ref{chap:ch16_climatology} does \emph{not} re-open), the
weak term merely bounds slow drift, and SST is now \emph{free to respond} to
$\Forc$ on the multi-year anomaly timescale. The same single function selects the
two modes by the presence or absence of \code{q\_flux}:

\begin{lstlisting}[caption={One function, two modes: control restoring vs free q-flux.}]
def heat_correction(sst_K, sst_target_K, restore_tau_days, q_flux=None):
    corr = restoring_flux(sst_K, sst_target_K, restore_tau_days)  # weak or strong
    if q_flux is not None:          # FREE mode: add frozen q-flux
        corr = corr + jnp.asarray(q_flux)
    return corr
\end{lstlisting}
\codref{chronos\_esm/ocean/flux\_correction.py}{\code{heat\_correction} ---
Equation~\eqref{eq:f14_free}; \code{q\_flux=None} recovers
Equation~\eqref{eq:f14_haney}}

\begin{remark}[Differentiability survives the correction]
Every operation in Equations~\eqref{eq:f14_myhre}--\eqref{eq:f14_free} is
expressed in \code{jnp}, with no \code{numpy} or \code{pint} hops. Consequently
$\mathrm{d}(\text{climate})/\mathrm{d}(\text{forcing})$ survives the flux
correction: the q-flux is a \emph{constant additive field}, so it contributes
nothing to the gradient with respect to $C$, while the weak restoring contributes
a bounded, differentiable damping. This is what makes the forced experiment a
legitimate target for the inverse methods of
Chapter~\ref{chap:ch15_differentiable_applications}.
\end{remark}

% ======================================================================
\section{The abrupt-2\texorpdfstring{$\times$}{x}CO\texorpdfstring{$_2$}{2} experiment}\label{sec:f14_experiment}
% ======================================================================

The experiment, following the spirit of the CMIP \texttt{abrupt-2xCO2} protocol
\citep{eyring2016cmip6}, is driven by \codref{experiments/run\_dino\_co2.py}{the
forced-CO\textsubscript{2} driver}. Its design has three deliberate features.

\paragraph{Branching from a common state.} Both branches start from the
\emph{same} converged control checkpoint and its frozen $\Qflux$. We integrate
two free-mode trajectories that differ only in the prescribed concentration:
\begin{align}
  \text{baseline:} \quad & C = C_0 = 280\,\mathrm{ppm}
    &&\Rightarrow\; \Forc = 0, \label{eq:f14_base}\\
  \text{forced:} \quad & C = 2C_0 = 560\,\mathrm{ppm}
    &&\Rightarrow\; \Forc = 5.35\ln 2 \approx 3.71\,\mathrm{W\,m^{-2}}.
    \label{eq:f14_forced}
\end{align}

\paragraph{Warming as a difference (drift cancellation).} The free model still
has a small residual drift (the weak restoring does not pin the mean exactly).
The CO\textsubscript{2} warming is defined as the \emph{difference} of the two
ocean-mean SST trajectories,
\begin{equation}
  \Delta\SST(t) \;=\; \SST^{\,\mathrm{forced}}(t) - \SST^{\,\mathrm{baseline}}(t),
  \label{eq:f14_dSST}
\end{equation}
because both branches inherit the \emph{same} drift from the same initial state
and the same $\Qflux$, so subtracting one from the other cancels it to leading
order and isolates the forced signal. The ocean-area-mean SST of each branch is

\begin{equation}
  \overline{\SST}(t) \;=\;
  \frac{\sum_{i\in\mathrm{ocean}} \SST^{(i)}(t)}{\sum_{i\in\mathrm{ocean}} 1}
  \;-\; 273.15 \quad [^\circ\mathrm{C}],
  \label{eq:f14_oceanmean}
\end{equation}
the sum running over the surface ocean grid cells (\code{omask}); see
\code{\_run\_branch} / \code{sst\_mean\_C} in the driver.

\paragraph{Time horizon.} The branches run for $\sim\!50$ model years
(\code{--years 50}, the experiment's default), stepping the fast coupled model at
the coupling interval and recording $\overline{\SST}$ annually. Fifty years is
short relative to the deep-ocean equilibration timescale ($\sim\!10^3$ yr), which
is exactly why the result is a transient, not an equilibrium, measure --- the
point of \S\ref{sec:f14_honesty}.

\begin{example}[Two lines, one number]
The driver prints the forcing $\Forc$ from Equation~\eqref{eq:f14_2x}, the
end-of-run warming $\Delta\SST(\text{end})$ from Equation~\eqref{eq:f14_dSST},
and their ratio $\Delta\SST/\Forc$ in $\mathrm{K\,(W\,m^{-2})^{-1}}$. That ratio
is the headline climate-response number, and it is what we interpret next.
\end{example}

% ======================================================================
\section{Result and its honest interpretation}\label{sec:f14_result}
% ======================================================================

A 50-year free-mode abrupt-$2\times$CO\textsubscript{2} run yields a global
ocean-mean surface warming of
\begin{equation}
  \Delta\SST \;\approx\; +1.58\,\mathrm{K},
  \qquad
  \frac{\Delta\SST}{\Forc} \;\approx\;
  \frac{1.58}{3.71} \;\approx\; 0.43\,\mathrm{K\,(W\,m^{-2})^{-1}}.
  \label{eq:f14_result}
\end{equation}
The model has moved, in the right direction, by a physically plausible amount:
the sign is correct, the magnitude is within the order of magnitude of observed
estimates, and (crucially) the response is \emph{nonzero}, confirming that the
q-flux successfully freed the SST from the restoring clamp of
Equation~\eqref{eq:f14_balance}.

% ----------------------------------------------------------------------
\subsection{What this number is --- and is not}\label{sec:f14_honesty}
% ----------------------------------------------------------------------

It would be easy, and wrong, to call $0.43\,\mathrm{K\,(W\,m^{-2})^{-1}}$ a
climate sensitivity. Three structural limitations make it a \emph{proxy} only,
and a textbook must state them plainly.

\begin{enumerate}[leftmargin=2em]
\item \textbf{Transient, not equilibrium.} Fifty years samples only the upper
  ocean's adjustment; the deep ocean continues to take up heat for centuries.
  $\Delta\SST$ in Equation~\eqref{eq:f14_result} is therefore a TCR-like
  \emph{transient} response, structurally smaller than the equilibrium climate
  sensitivity (ECS) the same forcing would eventually produce.

\item \textbf{No atmospheric radiative feedback amplification.} The forcing
  enters as a fixed surface-heat term (Equation~\eqref{eq:f14_inject}); there is
  \emph{no top-of-atmosphere energy closure}. The water-vapour, lapse-rate, cloud
  and surface-albedo feedbacks that amplify (or damp) the response in a full GCM
  are simply absent. The number is a \emph{surface-forcing} response, with the
  net feedback parameter effectively set by the weak restoring
  $\Hrest_{\mathrm{weak}}$ and the ocean's heat uptake rather than by radiative
  physics. The driver itself prints the caveat ``TCR-like proxy, NOT equilibrium
  ECS''.

\item \textbf{Cold-tropics base state.} The control climate carries a cold
  tropical bias (the cold-tongue / dead-ITCZ issue of
  Chapter~\ref{chap:ch16_climatology}). Surface fluxes --- the longwave
  $\propto \SST^4$ and the Clausius--Clapeyron-controlled latent flux --- are
  nonlinear in SST, so a response computed about a biased base state is itself
  biased; the slope in Equation~\eqref{eq:f14_result} should not be read to two
  significant figures of accuracy.
\end{enumerate}

\begin{proposition}[The result is a lower-bound-flavoured TCR proxy]
Because the dominant missing ingredients --- deep-ocean equilibration time and
positive water-vapour/cloud feedbacks --- both act to \emph{increase} the
eventual warming relative to Equation~\eqref{eq:f14_result}, while no comparably
large omitted process is known to decrease it, the model's
$0.43\,\mathrm{K\,(W\,m^{-2})^{-1}}$ is best read as a conservative,
transient-flavoured estimate. It is a demonstration that the coupled model
\emph{responds correctly in sign and order of magnitude}, not a calibrated
sensitivity.
\end{proposition}

\begin{remark}[Why this is still scientifically useful]
A differentiable model whose forced response is honest about its scope is
precisely the right laboratory for \emph{sensitivity analysis} and
\emph{calibration}. The very next chapter
(\ref{chap:ch15_differentiable_applications}) uses the differentiability
preserved through the q-flux to compute $\mathrm{d}\Delta\SST/\mathrm{d}C$ by
automatic differentiation rather than by finite differencing two 50-year runs ---
turning the experiment of this chapter into a single gradient evaluation.
\end{remark}

\begin{table}[ht]
\centering
\caption{Key quantities of the abrupt-$2\times$CO\textsubscript{2} experiment.}
\begin{tabular}{lll}
\toprule
Quantity & Value & Source \\
\midrule
Forcing coefficient & $5.35\,\mathrm{W\,m^{-2}}$ & \citet{myhre1998}, Eq.~\eqref{eq:f14_myhre} \\
$\Forc(2\times\text{CO}_2)$ & $3.71\,\mathrm{W\,m^{-2}}$ & Eq.~\eqref{eq:f14_2x} \\
Control restoring $\Hrest$ ($\tau{=}30\,$d) & $79\,\mathrm{W\,m^{-2}K^{-1}}$ & Eq.~\eqref{eq:f14_lambda30} \\
Free restoring $\Hrest_{\mathrm{weak}}$ ($\tau{=}3650\,$d) & $0.65\,\mathrm{W\,m^{-2}K^{-1}}$ & Eq.~\eqref{eq:f14_lambdaweak} \\
Warming $\Delta\SST$ (50 yr) & $+1.58\,\mathrm{K}$ & Eq.~\eqref{eq:f14_result} \\
Response $\Delta\SST/\Forc$ & $0.43\,\mathrm{K\,(W\,m^{-2})^{-1}}$ & Eq.~\eqref{eq:f14_result} \\
\bottomrule
\end{tabular}
\end{table}

% ======================================================================
\section{Adding 4\texorpdfstring{$\times$}{x}CO\texorpdfstring{$_2$}{2}: a sublinear response}
\label{sec:f14_4x}
% ======================================================================
Running the same protocol at $4\times$CO\textsubscript{2} ($1120$~ppm) gives a second point
on the forcing-response curve (Figure~\ref{fig:f14_response}). The Myhre forcing doubles,
$\Forc(4\times)=5.35\ln 4 = 7.42\,\mathrm{W\,m^{-2}}$, but the warming does \emph{not}:
\begin{equation}
  \Delta\SST(4\times) \approx +2.35\,\mathrm{K},\qquad
  \frac{2.35}{7.42}\approx 0.32\,\mathrm{K\,(W\,m^{-2})^{-1}}.
\end{equation}
The second doubling adds only $\approx 0.77\,\mathrm{K}$ against the first doubling's
$1.58\,\mathrm{K}$, so the response is \emph{sub-logarithmic} in CO\textsubscript{2}: the
warming per unit forcing falls from $0.43$ to $0.32\,\mathrm{K\,(W\,m^{-2})^{-1}}$. A truly
linear-in-forcing model would warm $\approx 3.16\,\mathrm{K}$ at $4\times$; the model warms
$0.8\,\mathrm{K}$ less. This is consistent with a proxy that lacks the amplifying atmospheric
radiative feedbacks (which in nature make the response closer to linear in $\ln$CO$_2$) and
with the surface-flux response saturating as the freed mixed layer warms. The end-of-run
$\Delta\SST$ carries $\sim\!\pm0.3\,\mathrm{K}$ of interannual noise, smaller than the
$0.8\,\mathrm{K}$ departure from linearity, so the curvature is a genuine model property,
not noise. Both points should still be read as transient TCR-like \emph{proxies}, not ECS.

\begin{figure}[ht]
\centering
\includegraphics[width=0.7\linewidth]{forcing_response}
\caption{Transient global-mean surface warming $\Delta\SST$ versus CO\textsubscript{2}
radiative forcing for the abrupt-$2\times$ and abrupt-$4\times$CO\textsubscript{2}
experiments. The dashed line is the linear response calibrated on the $2\times$ point; the
$4\times$ point falls below it, showing the sub-logarithmic (saturating) forcing response of
the proxy. Regenerate with \code{experiments/plot\_forcing\_response.py}.}
\label{fig:f14_response}
\end{figure}

% ======================================================================
\section*{Exercises}
% ======================================================================

\begin{exercise}[The forcing trap, quantified]
Using Equation~\eqref{eq:f14_balance}, compute the steady SST anomaly that the
\emph{control} restoring ($\Hrest=79\,\mathrm{W\,m^{-2}K^{-1}}$) would permit
under $2\times$CO\textsubscript{2}, and compare it with the
$0.65\,\mathrm{W\,m^{-2}K^{-1}}$ free mode. By what factor does freeing the ocean
increase the permitted steady response? Explain in one sentence why this factor
is \emph{not} the same as the ratio of the actual simulated warmings, and which
two physical processes (in the transient run) break the simple proportionality.
\end{exercise}

\begin{exercise}[A forcing sweep]
Run \code{experiments/run\_dino\_co2.py} for $C\in\{420,560,1120\}\,\mathrm{ppm}$
against the $280\,\mathrm{ppm}$ baseline. For each, record $\Forc$ from
Equation~\eqref{eq:f14_myhre} and the simulated $\Delta\SST$. Plot $\Delta\SST$
against $\Forc$ (not against $C$). Is the relationship linear? Fit a slope and
compare it with Equation~\eqref{eq:f14_result}. What does the logarithmic form of
Equation~\eqref{eq:f14_myhre} imply for the spacing of your three points on the
\emph{concentration} axis versus the \emph{forcing} axis?
\end{exercise}

\begin{exercise}[Drift cancellation matters]
The experiment defines warming as a difference of two branches
(Equation~\eqref{eq:f14_dSST}). Re-derive $\Delta\SST$ from the single forced
branch alone (do \emph{not} subtract the baseline) and compare. Estimate the
residual drift of the free model from the baseline curve and argue, with numbers,
whether ignoring it would have changed the headline
$0.43\,\mathrm{K\,(W\,m^{-2})^{-1}}$ by more than its honest uncertainty.
\end{exercise}

\begin{exercise}[From two runs to one gradient]
The Myhre forcing is differentiable in $C$. Sketch how you would obtain
$\mathrm{d}\Delta\SST/\mathrm{d}C$ at $C=560\,\mathrm{ppm}$ using \code{jax.grad}
through \code{ocean\_fluxes\_jax} rather than by finite-differencing two 50-year
integrations. Which term in the surface budget carries the gradient, and which
two terms (the frozen $\Qflux$ and the area-mean removal) contribute zero or a
constant to it? Cross-reference your answer with
Chapter~\ref{chap:ch15_differentiable_applications}.
\end{exercise}
